% Part: incompleteness
% Chapter: representability-in-q
% Section: c-representable

\documentclass[../../../include/open-logic-section]{subfiles}

\begin{document}

\olfileid{inc}{req}{cre}
\olsection{$C$નાં વિધેયોનું $\Th{Q}$માં પ્રતિનિધિત્વ કરી શકાય છે}

$C$માંના દરેક વિધેયનું
$\Th{Q}$માં પ્રતિનિધિત્વ
કરી શકાય છે તે
બતાવવાનું છે. અંતે,
$C$માંના દરેક $k$-સ્થાની
વિધેય $f(x_0,\dots,x_{k-1})$
માટે તેનું પ્રતિનિધિત્વ
કરતું સૂત્ર
$!A_f(x_0,\dots,x_{k-1},y)$
કેવી રીતે આપવું
તે બતાવવું પડશે.

\begin{lem}
પ્રાકૃતિક સંખ્યાઓ $n$
અને $m$ માટે, જો
$n \neq m$, તો
$Q \Proves \num
n \neq \num m$.
\end{lem}

\begin{proof}
$n$ પરના આગમનથી
બતાવો કે દરેક $m$
માટે, જો $n \neq m$,
તો $Q \Proves \eq/[\num n][\num m]$.

પાયાના કિસ્સામાં
$n = 0$. જો $m$
એ $0$ નથી, તો
કોઈ પ્રાકૃતિક
સંખ્યા $k$ માટે
$m = k + 1$.
આપણી પાસે
$\lforall[x][\eq/[0][x']]$
એવું સ્વયંસિદ્ધિ છે.
પરિમાણક-સ્વયંસિદ્ધિથી
$x$ની જગ્યાએ
$\num k$ મૂકતાં
$\eq/[0][\num k']$
મળે છે. પરંતુ
$\num k'$ એ
$\num m$ જ છે.

આગમન-પગલામાં
દાવો $n$ માટે સાચો
છે એમ માની $n+1$
વિચારો. $m$ કોઈ
પણ પ્રાકૃતિક
સંખ્યા હોય. બે
શક્યતાઓ છે:
કાં તો $m = 0$,
અથવા કોઈ $k$
માટે $m = k+1$.
પહેલો કિસ્સો
ઉપર મુજબ છે.
બીજા કિસ્સામાં
$n+1 \neq k+1$
માનો; ત્યારે
$n \neq k$.
$n$ માટેની
આગમન-પરિકલ્પનાથી
$\Th{Q} \Proves \eq/[\num n][\num k]$.
આપણી પાસે
$\lforall[x][\lforall[y][\eq[x'][y'] \lif \eq[x][y]]]$
સ્વયંસિદ્ધિ છે.
પરિમાણક-સ્વયંસિદ્ધિ
વાપરી
$\eq[\num n'][\num k'] \lif \eq[\num n][\num
  k]$ મળે. વિધાનાત્મક
તર્કથી $\Th{Q}$માં
$\eq/[\num n][\num k] \lif \eq/[\num n'][\num k']$
મળે છે. મોડસ
પોનેન્સથી ઇચ્છિત
$\eq/[\num n'][\num k']$
મળે છે, કારણ કે
$\num k'$ એ
$\num m$ છે.
\end{proof}

આ લેમ્મા બહુ મોટો
દાવો કરતું નથી:
મૂળભૂત રીતે, જુદા
સંખ્યાચિહ્નો જુદી
વસ્તુઓ નિર્દેશે છે
એ $\Th{Q}$ સાબિત
કરી શકે છે. દાખલા
તરીકે, $\Th{Q}$
$0'' \neq 0'''$
સાબિત કરે છે.
પરંતુ સામાન્ય કિસ્સો
બતાવવા થોડી
કાળજી જોઈએ.
આગમન વાપરીએ
છીએ, પણ તે
$\Th{Q}$ની
\emph{બહારનું}
આગમન છે.

શૂન્ય, ઉત્તરગામી વિધેય, સરવાળો, ગુણાકાર, સમાનતા માટેનું
લાક્ષણિક વિધેય અને પ્રક્ષેપોનું પ્રતિનિધિત્વ કરી શકાશે.
દરેક કિસ્સામાં પ્રતિનિધિત્વ કરતું યોગ્ય સૂત્ર સીધું છે.
દાખલા તરીકે, શૂન્યનું પ્રતિનિધિત્વ આ સૂત્ર કરે છે:
\[
y = 0,
\]
ઉત્તરગામી વિધેયનું આ સૂત્ર કરે છે:
\[
x_0' = y,
\]
અને સરવાળાનું આ સૂત્ર કરે છે:
\[
x_0 + x_1 = y.
\]
અહીં કામ એ સાબિત કરવાનું છે કે $\Th{Q}$ સંબંધિત વિધાનો
સાબિત કરી શકે છે. દાખલા તરીકે, ઉપરનું સૂત્ર સરવાળાનું
પ્રતિનિધિત્વ કરે છે એમ કહેવા માટે દરેક પ્રાકૃતિક સંખ્યા
$m$ અને $n$ માટે $\Th{Q}$ નીચેનાં બંને સૂત્રો સાબિત કરે છે
એ બતાવવું પડે:
\[
\eq[\num n + \num m][\num {n+m}]
\]
અને
\[
\lforall[y][(\eq[(\num n + \num m)][y] \lif \eq[y][\num {n+m}])].
\]

હવે સંયોજનનું શું? ધારો કે $h$ આમ વ્યાખ્યાયિત છે:
\[
h(x_0,\dots,x_{l-1}) = f(g_0(x_0,\dots,x_{l-1}), \dots,
g_{k-1}(x_0,\dots,x_{l-1})).
\]
અહીં અનુક્રમે $f,g_0,\dots,g_{k-1}$નું પ્રતિનિધિત્વ
કરતાં સૂત્રો $!A_f, !A_{g_0}, \dots,
!A_{g_{k-1}}$ આપણે પહેલેથી મેળવ્યાં છે. ત્યારે
$h$નું પ્રતિનિધિત્વ કરતું સૂત્ર $!A_h$ મેળવવા માટે
$!A_h(x_0,\dots,x_{l-1},y)$ને આમ વ્યાખ્યાયિત કરી શકીએ:
\begin{multline*}
  \lexists[z_0][\dots \lexists[z_{k-1}][(!A_{g_0}(x_0,\dots,x_{l-1},z_0) \land
  \dots \land !A_{g_{k-1}}(x_0,\dots,x_{l-1},z_{k-1}) \land\\
  !A_f(z_0,\dots,z_{k-1},y))]].
\end{multline*}
\sourcecorrection{OLMD-181}{મૂળ સૂત્રમાં અંતિમ $!A_f$ સંયોજક
અસ્તિત્વવાચક પરિમાણકોના વ્યાપની બહાર પડી જાય છે; અહીં બધા
મધ્યવર્તી $z_i$ અને અંતિમ સંયોજકને એ જ વ્યાપમાં રાખ્યાં છે.}

છેવટે અનિયત મર્યાદા સુધીની શોધ વિચારીએ. ધારો કે
$g(x,\vec z)$ નિયમિત છે અને $\Th{Q}$માં, માનો કે
$!A_g(x,\vec z,y)$ સૂત્રથી, તેનું પ્રતિનિધિત્વ કરી
શકાય છે. $f$ને $f(\vec z) = \mu x \; g(x,\vec z)$
વડે વ્યાખ્યાયિત કરો. $f$નું પ્રતિનિધિત્વ કરતું
$!A_f(\vec z,y)$ સૂત્ર શોધવું છે. સ્વાભાવિક પસંદગી આ છે:
\[
!A_f(\vec z,y) \ident !A_g(y,\vec z,0) \land \lforall[w][(w < y
\lif \lnot !A_g(w,\vec z,0))].
\]

નીચેના જેવાં વધારાનાં લેમ્મા વાપરી આ પસંદગી કામ કરે છે
એ બતાવી શકાય:

\begin{lem}
  દરેક ચલ $x$ અને દરેક પ્રાકૃતિક સંખ્યા $n$ માટે
  $\Th{Q}$ એ $\eq[(x'
  + \num n)][(x + \num n)']$ સાબિત કરે છે.
\end{lem}

ફરી નોંધવું જોઈએ કે આ દાવો $\Th{Q}$ એ
$\lforall[x][\lforall[y][(x' + y) = (x +y)']]$
સાબિત કરે છે એવા દાવા કરતાં નબળો છે; પછીનો દાવો ખોટો છે.

\begin{proof}
રાબેતા મુજબ સાબિતી $n$ પરના આગમનથી મળે છે.
પાયાના કિસ્સામાં $n = 0$ છે અને $\Th{Q}$ એ
$\eq[(x'+0)][(x+0)']$ સાબિત કરે છે એમ બતાવવાનું છે.
પરંતુ આપણી પાસે આ છે:
\begin{eqnarray*}
& & \eq[(x' + 0)][x'] \quad \text{સ્વયંસિદ્ધિ 4થી} \\
& & \eq[(x + 0)][x] \quad \text{સ્વયંસિદ્ધિ 4થી} \\
& & \eq[(x+0)'][x'] \quad \text{બીજી પંક્તિથી} \\
& & \eq[(x' + 0)][(x + 0)'] \quad \text{પહેલી અને ત્રીજી પંક્તિથી}
\end{eqnarray*}
આગમન-પગલામાં $\Th{Q}$માં
$\eq[(x' + \num n)][(x + \num n)']$ નિષ્પન્ન થયું
છે એમ ધારીએ. $\num{n+1}$ એ $\num n'$ હોવાથી,
$\Th{Q}$ એ $\eq[(x' + \num n')][(x + \num n')']$
સાબિત કરે છે એમ બતાવવાનું છે. $\Th{Q}$માં
સમાનતાઓની નીચેની સાંકળ નિષ્પન્ન કરી શકાય છે:
\begin{eqnarray*}
(x' + \num n') & \eq & (x' + \num n)' \quad \text{સ્વયંસિદ્ધિ 5થી} \\
& \eq & (x + \num n')' \quad \text{આગમન-પરિકલ્પનાથી}
\end{eqnarray*}
\sourcecorrection{OLMD-182}{બીજી સમાનતા માટે આગમન-પરિકલ્પના
ઉપરાંત $(x+\num n')=(x+\num n)'$ આપતું સ્વયંસિદ્ધિ 5
અને સમાનતાની અવેજી પણ જોઈએ; મૂળમાં આ મધ્યવર્તી
પગલાં સ્પષ્ટ લખ્યાં નથી.}
\end{proof}

\begin{lem}
\begin{enumerate}
\item $\Th{Q}$ એ $\lnot (x < \num 0)$ સાબિત કરે છે.
\item દરેક પ્રાકૃતિક સંખ્યા $n$ માટે $\Th{Q}$ એ
નીચેનું સાબિત કરે છે:
\[
x < \num {n+1} \lif (\eq[x][\num 0] \lor \dots \lor \eq[x][\num n]).
\]
\end{enumerate}
\end{lem}

\begin{proof}
પહેલો દાવો અને બીજા દાવાનો અમુક ભાગ અનૌપચારિક રીતે
સાબિત કરીએ; એટલે ઔપચારિક !!{derivation} કેવી રીતે
બનાવવી તેના માત્ર સંકેતો આપીશું.

પહેલા દાવા માટે $<$ની વ્યાખ્યા મુજબ $\Th{Q}$માં
$\lnot \lexists[y][\eq[(y'+x)][0]]$ સાબિત કરવાનું છે.
પ્રથમ-ક્રમ તર્કના સ્વયંસિદ્ધિઓ અને નિયમો વડે આ
$\lforall[y][\eq/[(y' + x)][0]]$ની સમકક્ષ છે.
વિચાર આ છે: $\eq[(y'+x)][0]$ ધારો. જો $x$ એ
શૂન્ય હોય, તો $\eq[(y' + 0)][0]$ મળે. પરંતુ
$\Th{Q}$ની સ્વયંસિદ્ધિ 4થી $\eq[(y' + 0)][y']$
મળે છે અને સ્વયંસિદ્ધિ 2થી $\eq/[y'][0]$
મળે છે; આ વિરોધાભાસ છે. જો $x$ એ $0$
ન હોય, તો સ્વયંસિદ્ધિ 3થી કોઈ $z$ માટે
$\eq[x][z']$ મળે છે. ત્યારે
$\eq[(y' + z')][0]$ મળે. સ્વયંસિદ્ધિ~5થી
$\eq[(y' + z)'][0]$ મળે, જે ફરી
સ્વયંસિદ્ધિ~2નો વિરોધ કરે છે. બંને કિસ્સાઓ
પૂરા થતાં $\lforall[y][\eq/[(y' +x)][0]]$ મળે છે.
\sourcecorrection{OLMD-183}{મૂળ ગદ્યમાં બંને
કિસ્સાઓ પૂરા થાય તે પહેલાં જ સર્વવાચી નિષ્કર્ષ
લખાયો છે; અહીં તેને બીજા કિસ્સા પછી મૂક્યો છે.}

બીજા દાવા માટે $n$ પર આગમન વાપરો. પાયાનો
કિસ્સો $n =0$ છે. ત્યારે
$x < \num 1 \lif \eq[x][\num 0]$ બતાવવાનું છે.
$x < \num 1$ ધારો. $<$ માટેની વ્યાખ્યાત્મક
સ્વયંસિદ્ધિથી $\lexists[y][\eq[(y'+x)][0']]$
મળે છે. એવો ગુણધર્મ ધરાવતું $y$ લો;
તેથી $\eq[y'+x][0']$ મળે છે.

$\eq[x][0]$ બતાવવાનું છે. સ્વયંસિદ્ધિ 3થી,
જો $x$ એ શૂન્ય ન હોય, તો કોઈ $z$ માટે તે
$z'$ને સમાન છે. તેથી $\eq[(y' + z')][0']$
મળે છે. $\Th{Q}$ની સ્વયંસિદ્ધિ~5થી
$\eq[(y'+z)'][0']$ મળે. સ્વયંસિદ્ધિ 1થી
$\eq[(y' + z)][0]$ મળે. પરંતુ વ્યાખ્યા મુજબ
તેનો અર્થ $z < 0$ થાય છે, જે પહેલા દાવાનો
વિરોધ કરે છે.
\sourcecorrection{OLMD-184}{મૂળ પાયાના કિસ્સામાં
$x<\num 1$થી $x=0$ દિશા જ સમજાવી છે. વિપરીત
દિશા $x=0$ હોય ત્યારે $<$ની વ્યાખ્યામાં
$y=0$ સાક્ષી લઈને મળે છે; આગળનું આગમન-પગલું
પણ આ જૂના સંક્ષિપ્ત વિભાગમાં લખાયું નથી.}
\end{proof}

\begin{explain}
આપણે બતાવ્યું છે કે સંગણનીય વિધેયોનો ગણ
$\Th{Q}$માં પ્રતિનિધિત્વ કરી શકાય એવાં
વિધેયોના ગણ તરીકે ઓળખી શકાય છે. હકીકતમાં
સાબિતી વધુ સામાન્ય છે. પ્રતિનિધિત્વની
વ્યાખ્યાથી સમજાય છે કે $\Th{Q}$નું વિસ્તરણ
હોય એવી કોઈ પણ થિયરી, અથવા જેમાં $\Th{Q}$નું
અર્થઘટન કરી શકાય એવી થિયરી, સંગણનીય
વિધેયોનું પ્રતિનિધિત્વ કરી શકે છે. બીજી
બાજુ, સાબિતીની સંકલ્પના સંગણનીય હોય એવી
કોઈ પણ !!{derivation} પ્રણાલીમાં દરેક
પ્રતિનિધિત્વ કરી શકાય એવું વિધેય
સંગણનીય છે. તેથી સંગણનીય વિધેયોના
ગણને પેયાનો અંકગણિતમાં, અથવા ઝેર્મેલો--ફ્રેન્કેલ
ગણસિદ્ધાંતમાં પણ, પ્રતિનિધિત્વ થતાં
વિધેયોના ગણ તરીકે ઓળખી શકાય છે.
ગેડેલે નોંધ્યું તેમ, આ કંઈક આશ્ચર્યજનક છે.
સાબિત કરી શકાય તેવા વિધાનો વિશેના પ્રશ્નો
પસંદ કરેલી થિયરી પર ખૂબ આધાર રાખે છે
એ આપણે જોઈશું: સામાન્ય રીતે સ્વયંસિદ્ધિઓ
વધુ પ્રબળ હોય તેમ વધુ સાબિત કરી શકાય.
પરંતુ સ્વયંસિદ્ધિ-આધારિત થિયરીઓના
વિશાળ વર્ગમાં પ્રતિનિધિત્વ કરી શકાય
એવાં વિધેયો ચોક્કસ સંગણનીય વિધેયો જ છે.
\end{explain}

\end{document}
